\documentclass[10pt,twocolumn]{article}

\usepackage[margin=1in]{geometry}
\usepackage{amsmath,amssymb,amsfonts}
\usepackage{booktabs}
\usepackage{graphicx}
\usepackage{hyperref}
\usepackage{xcolor}
\usepackage{multirow}
\usepackage{algorithm}
\usepackage{algorithmic}
\usepackage{times}

\title{ACR-VLA: Adaptive Conformal Refinement for\\Vision-Language-Action Models}

\author{Anonymous Authors}
\date{CoRL 2026 Submission}

\begin{document}
\maketitle

%==============================================================================
\begin{abstract}
%==============================================================================
Flow-matching Vision-Language-Action (VLA) models achieve near-saturation performance on standard manipulation benchmarks but suffer catastrophic degradation under distribution shift.
We propose ACR-VLA, an inference-time framework that augments any flow-matching VLA with structured uncertainty estimation and learned action correction.
ACR-VLA consists of three modules:
(1)~\textbf{Conformal Generative Scoring (CGS)}, which draws $K$ flow-matching samples and applies multivariate conformal prediction with Ledoit-Wolf shrinkage to produce 45-dimensional per-DOF uncertainty features;
(2)~\textbf{Adaptive Conformal Inference (ACI)} with a novel \textbf{Chunk Overlap Consistency (COC)} signal that adapts conformal thresholds online without external supervision;
and (3)~a lightweight \textbf{Conformal Refinement Module (CRM)} (131K parameters) that takes base actions, visual embeddings, and conformal features as input to predict corrective residuals.
On LIBERO-Plus, ACR-VLA improves SmolVLA success rate by $+6$ percentage points (62\%$\to$68\%), with gains of $+9$--$10$\,pp on camera and noise perturbations.
Ablations confirm that conformal features provide a complementary conditioning signal beyond observation-only residual learning ($+3.5$\,pp, A3) and that CRM's learned correction outperforms sample selection ($+4$\,pp, A2).
When stacked on SRPO (82.1\%), ACR-VLA achieves 84.0\%, demonstrating architectural compatibility with training-time methods.
\end{abstract}

%==============================================================================
\section{Introduction}
\label{sec:intro}
%==============================================================================

Flow-matching Vision-Language-Action (VLA) models~have achieved remarkable success on standard manipulation benchmarks, with recent methods reaching 95--99\% success rates on LIBERO.
However, the LIBERO-Plus benchmark reveals a different picture: systematic perturbations in camera pose, observation noise, robot kinematics, lighting, backgrounds, scene layout, and language instructions expose dramatic robustness gaps, with success rates dropping to 50--65\% for state-of-the-art VLAs.

Existing approaches address this robustness challenge from different angles.
\emph{Training-time methods} such as SRPO, RobustVLA, and DiG-Flow modify the policy's learned representations through data augmentation, adversarial training, or diffusion guidance, achieving strong gains (e.g., SRPO: 82.1\% on LIBERO-Plus) but requiring retraining.
\emph{Inference-time methods} such as SAFE detect failure and trigger recovery, while ReconVLA uses scalar Euclidean distance for sample selection among multiple flow-matching candidates.
CP-in-the-Loop integrates conformal prediction into planning loops, and UPS uses conformal scores for policy gating.
However, no existing method \emph{directly conditions a learned correction module on structured, per-DOF conformal uncertainty features} from multivariate conformal prediction.

We identify a key gap: scalar uncertainty summaries (as in ReconVLA) discard the rich per-DOF structure of action uncertainty---which dimensions are uncertain and how they co-vary.
This information is precisely what a learned correction module needs to allocate its residual predictions effectively.
Our core insight is that multivariate conformal prediction on flow-matching samples produces structured uncertainty features---per-DOF standard deviations, conformal radii, and ellipsoid volume---that serve as a complementary conditioning signal for learned action correction.

We propose \textbf{ACR-VLA} (Adaptive Conformal Refinement for VLAs), a modular inference-time framework with three components:
\begin{enumerate}
    \item \textbf{Conformal Generative Scoring (CGS)}: Draws $K{=}10$ flow-matching samples, estimates the action distribution covariance via Ledoit-Wolf shrinkage, and computes multivariate conformal prediction sets yielding 45-dimensional uncertainty features (7 per-DOF standard deviations $+$ 7 conformal radii $+$ 21 unique covariance entries $+$ 7 per-DOF Mahalanobis contributions $+$ 3 summary statistics).
    \item \textbf{Adaptive Conformal Inference (ACI) with Chunk Overlap Consistency (COC)}: COC provides a self-referential coverage signal by measuring agreement between overlapping action chunks, enabling ACI to adapt conformal thresholds online without any external supervision or reward signal.
    \item \textbf{Conformal Refinement Module (CRM)}: A lightweight 131K-parameter network that takes the base VLA action, visual embeddings, proprioceptive state, and CGS-derived conformal features as input to predict a corrective residual.
\end{enumerate}

Our contributions are:
\textbf{(1)}~The CGS+CRM framework, the first method to use structured multivariate conformal features as direct input to a learned action correction module.
\textbf{(2)}~COC, a supervision-free proxy signal enabling online ACI adaptation during deployment without environment reward access.
\textbf{(3)}~Comprehensive experiments on LIBERO-Plus demonstrating $+6$\,pp improvement on SmolVLA and architectural compatibility with SRPO, alongside detailed ablations and feature importance analysis.

%==============================================================================
\section{Related Work}
\label{sec:related}
%==============================================================================

\paragraph{Robust VLA Policies.}
Training-time approaches modify the policy itself.
SRPO applies self-reinforced policy optimization to achieve 82.1\% on LIBERO-Plus.
RobustVLA uses data augmentation with diverse perturbations during fine-tuning.
DiG-Flow introduces diffusion guidance to steer flow-matching generation toward robust behaviors.
These methods are complementary to inference-time correction, as we demonstrate empirically by stacking ACR-VLA on SRPO.

\paragraph{Uncertainty-Aware Action Selection.}
ReconVLA samples multiple flow-matching candidates and selects the one with lowest Euclidean reconstruction distance---a scalar summary that discards per-DOF structure.
SAFE monitors uncertainty to detect and recover from failures.
UPS uses conformal prediction scores for policy gating (execute vs.\ abstain) but does not use conformal features as correction network inputs.
In contrast, ACR-VLA feeds structured conformal features into a learned module that synthesizes corrected actions rather than merely selecting or gating.

\paragraph{Conformal Prediction in Robotics.}
Conformal prediction provides distribution-free coverage guarantees.
CP-in-the-Loop integrates conformal sets into planning loops for trajectory optimization.
Multivariate extensions using Mahalanobis distance enable joint prediction regions that capture cross-DOF correlations.
Our CGS module builds on this foundation, using Ledoit-Wolf shrinkage to stabilize covariance estimation from $K{=}10$ samples, and novel application of ACI with a self-referential COC signal.

\paragraph{Learned Residual Policies.}
A2C2 and PLD learn residual corrections on top of base policies but do not incorporate calibrated uncertainty features.
ACR-VLA's CRM differs by conditioning the residual prediction on structured conformal statistics, providing the correction module with explicit information about which dimensions and how much the base policy is uncertain.

%==============================================================================
\section{Method}
\label{sec:method}
%==============================================================================

ACR-VLA operates as an inference-time wrapper around any flow-matching VLA. Given a base VLA policy $\pi_\theta$, at each time step $t$:

\subsection{Conformal Generative Scoring (CGS)}
\label{sec:cgs}

We draw $K{=}10$ independent flow-matching samples $\{\mathbf{a}^{(k)}\}_{k=1}^K$ from $\pi_\theta(\cdot | \mathbf{o}_t)$, where each $\mathbf{a}^{(k)} \in \mathbb{R}^7$ (7-DOF robot action).

\paragraph{Covariance Estimation.}
With $K{=}10$ and $p{=}7$, the sample-to-dimension ratio $K/p \approx 1.4$ makes the sample covariance singular or ill-conditioned.
We apply Ledoit-Wolf shrinkage:
\begin{equation}
\hat{\boldsymbol{\Sigma}}_{\mathrm{LW}} = (1 - \alpha_{\mathrm{LW}}) \hat{\boldsymbol{\Sigma}}_{\mathrm{sample}} + \alpha_{\mathrm{LW}} \cdot \mathrm{tr}(\hat{\boldsymbol{\Sigma}}_{\mathrm{sample}})/p \cdot \mathbf{I}_p,
\end{equation}
where $\alpha_{\mathrm{LW}}$ is the analytically optimal shrinkage intensity.

\paragraph{Nonconformity Score.}
We compute the Mahalanobis distance for each sample:
\begin{equation}
s^{(k)} = (\mathbf{a}^{(k)} - \bar{\mathbf{a}})^\top \hat{\boldsymbol{\Sigma}}_{\mathrm{LW}}^{-1} (\mathbf{a}^{(k)} - \bar{\mathbf{a}}),
\end{equation}
where $\bar{\mathbf{a}} = \frac{1}{K}\sum_k \mathbf{a}^{(k)}$.
The conformal prediction set at level $1{-}\alpha$ is defined by threshold $q_\alpha$: the $\lceil(1{-}\alpha)(K{+}1)\rceil/K$ quantile of calibration nonconformity scores.

\paragraph{Uncertainty Feature Extraction.}
From the $K$ samples and shrunk covariance, we extract a 45-dimensional feature vector $\mathbf{f}_{\mathrm{conf}} \in \mathbb{R}^{45}$:
\begin{itemize}
    \item Per-DOF standard deviations: $\sigma_j = \sqrt{[\hat{\boldsymbol{\Sigma}}_{\mathrm{LW}}]_{jj}}$ for $j=1,\ldots,7$
    \item Per-DOF conformal radii: $r_j = \sigma_j \sqrt{q_\alpha}$ for $j=1,\ldots,7$
    \item Upper-triangular covariance entries: 21 unique elements of $\hat{\boldsymbol{\Sigma}}_{\mathrm{LW}}$
    \item Per-DOF Mahalanobis contributions: 7 diagonal elements of the per-sample Mahalanobis decomposition
    \item Summary statistics: ellipsoid volume $V = \det(\hat{\boldsymbol{\Sigma}}_{\mathrm{LW}})^{1/2} \cdot q_\alpha^{p/2} \cdot V_p$, trace, and condition number (3 features)
\end{itemize}

\subsection{Adaptive Conformal Inference with Chunk Overlap Consistency}
\label{sec:aci}

Standard conformal prediction assumes exchangeability between calibration and test data. Under distribution shift, the fixed threshold $q_\alpha$ produces miscalibrated prediction sets. ACI addresses this by online updating:
\begin{equation}
\alpha_{t+1} = \alpha_t + \gamma(\alpha_{\mathrm{target}} - c_t),
\end{equation}
where $c_t \in \{0,1\}$ indicates whether the observation at step $t$ fell within the prediction set, and $\gamma$ is the step size.

\paragraph{Chunk Overlap Consistency (COC).}
In action-chunking VLAs, consecutive predictions produce overlapping action sequences. COC measures agreement in the overlapping region:
\begin{equation}
\mathrm{COC}_t = 1 - \frac{\|\mathbf{a}_{t}^{\mathrm{chunk}_{t-1}} - \mathbf{a}_{t}^{\mathrm{chunk}_{t}}\|}{\tau_{\mathrm{COC}}},
\end{equation}
where $\mathbf{a}_{t}^{\mathrm{chunk}_{t-1}}$ is the action at step $t$ predicted by the previous chunk, and $\tau_{\mathrm{COC}}$ is a normalization threshold.
High COC indicates the policy is temporally consistent; low COC suggests distribution shift or policy confusion.
COC serves as the coverage signal $c_t$ for ACI, enabling online threshold adaptation without any external supervision.

COC achieves Spearman $\rho{=}0.40$ ($p{=}0.002$) correlation with episode-level task success---a weak but statistically significant proxy.
Per-step binary classification accuracy is 62\%, sufficient for the slow exponential-moving-average dynamics of ACI (convergence in ${\sim}120$ episodes).

\subsection{Conformal Refinement Module (CRM)}
\label{sec:crm}

CRM is a lightweight feedforward network (131K parameters) that predicts a corrective residual:
\begin{equation}
\Delta\mathbf{a}_t = \mathrm{CRM}([\bar{\mathbf{a}}_t; \mathbf{e}_t^{\mathrm{vis}}; \mathbf{p}_t; \mathbf{f}_{\mathrm{conf},t}]),
\end{equation}
where $\bar{\mathbf{a}}_t$ is the mean base action, $\mathbf{e}_t^{\mathrm{vis}}$ is the visual embedding from the VLA encoder, $\mathbf{p}_t$ is proprioception, and $\mathbf{f}_{\mathrm{conf},t}$ is the 45D conformal feature vector.
The corrected action is $\mathbf{a}_t^* = \bar{\mathbf{a}}_t + \Delta\mathbf{a}_t$.

CRM is trained on LIBERO demonstration data with synthetic augmentation covering camera jitter, Gaussian action noise, and robot-state perturbation. The training loss is:
\begin{equation}
\mathcal{L}_{\mathrm{CRM}} = \mathbb{E}\left[\|\mathbf{a}_t^{\mathrm{demo}} - (\bar{\mathbf{a}}_t + \mathrm{CRM}(\cdot))\|^2\right].
\end{equation}

The architecture consists of two hidden layers with ReLU activations and layer normalization, deliberately kept lightweight to prevent overfitting and maintain inference speed (1.8\,ms per forward pass).

%==============================================================================
\section{Experiments}
\label{sec:experiments}
%==============================================================================

\subsection{Setup}

\paragraph{Base VLA.} We use SmolVLA, a flow-matching VLA achieving 95.0\% on LIBERO (10-task nominal suite).

\paragraph{Benchmarks.}
\emph{LIBERO} (nominal): 10 standard manipulation tasks, evaluating that ACR-VLA does not degrade in-distribution performance.
\emph{LIBERO-Plus}: Extends LIBERO with 7 perturbation types at 3 severity levels (L3--L5), testing robustness under distribution shift. We report per-level and average success rates.

\paragraph{Baselines and Ablations.}
We compare ACR-VLA against the base SmolVLA and SRPO (82.1\% LIBERO-Plus SOTA), and evaluate the following ablations:
\textbf{A1}: Static calibration (no ACI);
\textbf{A2}: Sample selection only (no CRM);
\textbf{A3}: CRM without conformal features;
\textbf{A4}: CRM with scalar uncertainty;
\textbf{A5}: Bonferroni per-DOF intervals vs.\ joint Mahalanobis;
\textbf{A6}: Oracle episode success vs.\ COC for ACI;
\textbf{A7}: DiU aleatoric vs.\ CGS epistemic features.

\subsection{Main Results}

Table~\ref{tab:main} presents the primary results.
ACR-VLA improves SmolVLA by $+6.0$\,pp on LIBERO-Plus (62.0\%$\to$68.0\%), with consistent gains across all severity levels ($+6$\,pp each at L3, L4, L5).
The nominal LIBERO degradation is 2.0\,pp (95.0\%$\to$93.0\%), arising from occasional over-correction on in-distribution steps where conformal features show mild uncertainty.

\begin{table}[t]
\centering
\small
\caption{Main results on LIBERO and LIBERO-Plus. SR = success rate (\%).}
\label{tab:main}
\begin{tabular}{@{}lccccc@{}}
\toprule
Method & LIB & L3 & L4 & L5 & LP-Avg \\
\midrule
SmolVLA & 95.0 & 72.0 & 60.0 & 54.0 & 62.0 \\
\;+ ACR-VLA & 93.0 & 78.0 & 66.0 & 60.0 & 68.0 \\
\;$\Delta$ & $-$2.0 & +6.0 & +6.0 & +6.0 & +6.0 \\
\midrule
SRPO & 97.0 & 87.0 & 81.0 & 78.3 & 82.1 \\
\;+ ACR-VLA & 96.5 & 89.0 & 83.0 & 80.0 & 84.0 \\
\;$\Delta$ & $-$0.5 & +2.0 & +2.0 & +1.7 & +1.9 \\
\bottomrule
\end{tabular}
\end{table}

When stacked on SRPO, ACR-VLA provides $+1.9$\,pp (82.1\%$\to$84.0\%), with minimal nominal degradation ($-0.5$\,pp).
This gain is modest and below the 2\,pp threshold we set for strong complementarity evidence.
We interpret this as \emph{complementary in principle}: SRPO's training-time optimization captures most correctable error, and the residual gains from inference-time correction are diminished but present.
The result demonstrates architectural compatibility---CRM can wrap any flow-matching VLA without interference---rather than a compelling empirical lift beyond training-time methods.

\paragraph{Inference Overhead.}
CGS adds 12.5\,ms (10 flow-matching forward passes) and CRM adds 1.8\,ms, totaling ${\sim}8\%$ pipeline overhead.
With optional SnapFlow 1-step distillation, CGS time drops to 18\,ms total (3.2$\times$ speedup) at a cost of 1\,pp SR.

\subsection{Ablation Study}

Table~\ref{tab:ablation} summarizes the ablation results.

\begin{table}[t]
\centering
\small
\caption{Ablation study on LIBERO-Plus. $\Delta$ vs.\ full ACR-VLA.}
\label{tab:ablation}
\begin{tabular}{@{}lccr@{}}
\toprule
Variant & LIB & LP & $\Delta$ \\
\midrule
ACR-VLA (full) & 93.0 & 68.0 & --- \\
\midrule
A1: No ACI (static) & 93.0 & 66.5 & $-$1.5 \\
A2: Selection only & 95.5 & 64.0 & $-$4.0 \\
A3: No conformal feat. & 94.0 & 64.5 & $-$3.5 \\
A4: Scalar uncertainty & 93.5 & 66.0 & $-$2.0 \\
A5: Bonferroni & --- & 66.5 & $-$1.5 \\
A7: DiU aleatoric & 94.5 & 65.0 & $-$3.0 \\
A7: Combined (52D) & 92.5 & 68.5 & +0.5 \\
\midrule
SmolVLA (base) & 95.0 & 62.0 & $-$6.0 \\
\bottomrule
\end{tabular}
\end{table}

\paragraph{A3: Conformal features as complementary signal.}
Removing conformal features from CRM (A3) reduces LP-Avg by 3.5\,pp (68.0\%$\to$64.5\%).
The A3 variant is equivalent to an observation-conditioned residual learner (similar to A2C2) and still achieves $+2.5$\,pp over the base VLA, confirming that learned residual correction has value independent of uncertainty conditioning.
The additional 3.5\,pp from conformal features demonstrates that CGS-derived uncertainty statistics provide a complementary conditioning signal for CRM beyond what observation features alone supply.
This gap clears our 2\,pp significance threshold but is moderate, consistent with conformal features serving as auxiliary information rather than the primary correction driver.

\paragraph{A2: Learned correction vs.\ sample selection.}
Sample selection (A2) achieves only $+2.0$\,pp over base, compared to CRM's $+6.0$\,pp.
The 4\,pp gap demonstrates that CRM's ability to \emph{synthesize} corrected actions---rather than merely select among existing samples---is the primary source of improvement.
When all $K$ samples are systematically biased under perturbation, selection cannot help; learned correction can.

\paragraph{A4: Structured vs.\ scalar uncertainty.}
Replacing 45D per-DOF features with a scalar uncertainty summary (A4) reduces LP-Avg by 2.0\,pp.
The ablation ladder---no uncertainty ($+2.5$\,pp), scalar ($+4.0$\,pp), structured 45D ($+6.0$\,pp) over base---shows that uncertainty features improve CRM regardless of representation, with structured features providing incremental value through per-DOF granularity.

\paragraph{A7: Epistemic vs.\ aleatoric uncertainty.}
CGS epistemic features outperform DiU aleatoric features by 3.0\,pp on LIBERO-Plus (68.0\% vs.\ 65.0\%).
Under distribution shift, epistemic uncertainty rises as the model encounters unfamiliar observations, providing a stronger signal for CRM than aleatoric variance, which is present even in-distribution.
Combining both (52D) yields only $+0.5$\,pp over epistemic-only, suggesting the aleatoric signal is largely subsumed.

\paragraph{A1: ACI contribution.}
ACI with COC provides $+1.5$\,pp over static calibration.
The modest gap reflects COC's noisy proxy nature ($\rho{=}0.40$); ACI's value is more in the principled online recalibration guarantee than in raw empirical lift.

\subsection{COC Signal Analysis}

\begin{table}[t]
\centering
\small
\caption{COC signal quality and ACI adaptation metrics.}
\label{tab:coc}
\begin{tabular}{lc}
\toprule
Metric & Value \\
\midrule
Spearman $\rho$ (COC, task success) & 0.40 ($p{=}.002$) \\
COC per-step binary accuracy & 62\% \\
ACI convergence & ${\sim}$120 episodes \\
ACI target / achieved coverage & 90\% / 88$\pm$4\% \\
Ellipsoid volume narrowing (vs.\ static) & 15\% \\
\midrule
COC-based ACI LP-Avg & 68.0\% \\
Oracle (episode success) ACI LP-Avg & 69.5\% \\
Supervision-freedom cost & $-$1.5\,pp \\
\bottomrule
\end{tabular}
\end{table}

Table~\ref{tab:coc} characterizes COC as an ACI signal.
COC's per-step accuracy (62\%) is limited---it measures temporal consistency of the policy, which correlates with but does not cause success.
However, the slow EMA dynamics of ACI are tolerant of per-step noise: the statistically significant episode-level correlation ($\rho{=}0.40$) is sufficient for ACI to converge within ${\sim}120$ episodes and achieve 15\% ellipsoid volume narrowing compared to static calibration.

The comparison with oracle episode-level success (A6) quantifies the \emph{price of supervision-freedom}: 1.5\,pp in SR and 3$\times$ slower convergence.
We argue this is an acceptable trade-off, as COC enables ACI adaptation during deployment without any external reward signal---a property essential for autonomous operation.

\subsection{Per-Perturbation Analysis}

\begin{table}[t]
\centering
\small
\caption{Per-perturbation ACR-VLA gain (pp) on LIBERO-Plus.}
\label{tab:perturb}
\begin{tabular}{@{}lcccc@{}}
\toprule
Perturbation & Base & Ours & $\Delta$ & Aug? \\
\midrule
Camera jitter & 58.0 & 67.0 & +9.0 & \checkmark \\
Obs.\ noise & 55.0 & 65.0 & +10.0 & \checkmark \\
Robot-state & 60.0 & 67.0 & +7.0 & part. \\
Lighting & 65.0 & 70.0 & +5.0 & --- \\
Background & 68.0 & 72.0 & +4.0 & --- \\
Layout & 62.0 & 64.0 & +2.0 & --- \\
Language & 66.0 & 67.0 & +1.0 & --- \\
\midrule
Avg (aug.) & --- & --- & +8.7 & \\
Avg (non-aug.) & --- & --- & +2.3 & \\
\bottomrule
\end{tabular}
\end{table}

Table~\ref{tab:perturb} reveals that ACR-VLA's gains concentrate on perturbation types covered by synthetic augmentation during CRM training.
Camera jitter ($+9$\,pp) and observation noise ($+10$\,pp) show the strongest improvements, while layout ($+2$\,pp) and language ($+1$\,pp) changes---which alter task semantics rather than inducing action-space distributional shift---show minimal benefit.
This is an expected and honest limitation: ACR-VLA corrects actions when the policy's uncertainty is due to observation-level perturbation, not when the task itself has changed.
For semantic perturbations, training-time methods that modify learned representations (SRPO, RobustVLA) remain more appropriate.

\subsection{Feature Importance Analysis}

\begin{table}[t]
\centering
\small
\caption{CRM feature importance (top 10 / 45 features).}
\label{tab:feature}
\begin{tabular}{@{}clcc@{}}
\toprule
\# & Feature & Grad. & Type \\
\midrule
1 & std(j1) & .180 & Marg. \\
2 & std(j4) & .150 & Marg. \\
3 & std(j6) & .140 & Marg. \\
4 & ellips.\ vol. & .120 & Summ. \\
5 & radius(j1) & .090 & Marg. \\
6 & radius(j4) & .075 & Marg. \\
7 & std(j5) & .060 & Marg. \\
8 & cov(j1,j4) & .045 & Cov. \\
9 & radius(j6) & .040 & Marg. \\
10 & cov(j4,j6) & .035 & Cov. \\
\bottomrule
\end{tabular}
\end{table}

Table~\ref{tab:feature} shows that CRM relies primarily on per-DOF marginal statistics (72\% of total gradient magnitude), with covariance-derived features contributing modestly (18\%) and the ellipsoid volume summary accounting for 10\%.
The dimension selectivity index is 0.35 (on a 0--1 scale), indicating that CRM pays differential attention to specific DOFs (proximal joints 1, 4, 6) but does not exhibit a clean ``correct only the uncertain DOFs'' pattern.
No perturbation-specific feature usage pattern was detected---CRM applies the same feature weighting across perturbation types.

This analysis motivates our reframing: the 45D structured features provide value primarily through richer marginal statistics (per-DOF std, radius) rather than through exploitation of cross-DOF covariance structure.
The multivariate CP formulation is justified by the marginal information it provides per-DOF, not by the covariance structure per se, though the joint Mahalanobis formulation does yield a modest 1.5\,pp advantage over Bonferroni (A5).

\subsection{Practical Considerations: SnapFlow Integration}

For deployment scenarios requiring real-time control, SnapFlow 1-step distilled generation reduces CGS inference from 58\,ms to 18\,ms (3.2$\times$ speedup) at a cost of 1\,pp SR (68.0\%$\to$67.0\%).
The KL divergence between SnapFlow and standard sample distributions is 0.08, with CGS ellipsoid volume 8\% smaller due to reduced tail variance in 1-step samples.
This enables ${\sim}55$\,Hz total pipeline throughput, making ACR-VLA feasible for real-time robotic control.

%==============================================================================
\section{Discussion}
\label{sec:discussion}
%==============================================================================

\paragraph{What CRM learns.}
The ablation decomposition reveals an instructive picture.
CRM without any uncertainty features (A3) achieves $+2.5$\,pp---the observation-conditioned residual already captures useful correction signal.
Adding conformal features provides an additional $+3.5$\,pp, confirming their value as complementary conditioning.
However, feature importance analysis shows CRM uses these features primarily through per-DOF marginals rather than covariance structure, and no perturbation-specific feature pattern emerges.
We conclude that structured uncertainty features from multivariate CP provide a complementary conditioning signal for learned action correction, but the mechanism is richer marginal information rather than dimension-selective correction in the strong sense.

\paragraph{COC as a deployment-ready signal.}
COC is not a reliable per-step success predictor (62\% accuracy), but it is a sufficient signal for the specific task of slow ACI threshold adaptation.
The 1.5\,pp price of supervision-freedom (vs.\ oracle episode success, A6) is a favorable trade-off for autonomous deployment where no reward signal is available.

\paragraph{Complementarity with training-time methods.}
The 1.9\,pp stacking gain on SRPO demonstrates architectural compatibility but also reveals diminishing returns: training-time optimization captures most low-hanging fruit, and inference-time correction provides marginal additional value.
We position ACR-VLA as \emph{complementary in principle} rather than claiming strong orthogonality.

\paragraph{Limitations.}
(1)~Gains on semantic perturbations (layout, language) are minimal ($+1$--$2$\,pp), as CGS detects action-space uncertainty rather than semantic task changes.
(2)~$K{=}10$ covariance estimation is noisy; the Ledoit-Wolf shrinkage toward identity attenuates the cross-DOF correlation structure that would enable stronger dimension-selective correction.
(3)~The 2\,pp nominal degradation, while acceptable, indicates CRM sometimes over-corrects in-distribution.
A learned gating mechanism could mitigate this.
(4)~All experiments are in simulation (LIBERO/LIBERO-Plus); real-world transfer remains to be validated.

%==============================================================================
\section{Conclusion}
\label{sec:conclusion}
%==============================================================================

We presented ACR-VLA, an inference-time framework that augments flow-matching VLAs with structured conformal uncertainty estimation and learned action correction.
The CGS+CRM architecture provides $+6$\,pp on LIBERO-Plus over SmolVLA, with ablations confirming that conformal features contribute a $+3.5$\,pp complementary conditioning signal beyond observation-only residual learning.
COC enables supervision-free online ACI adaptation at a modest 1.5\,pp cost compared to oracle signals.
ACR-VLA is architecturally compatible with training-time methods, achieving 84.0\% when stacked on SRPO.
The framework's modular design---CGS, ACI/COC, and CRM as independent components---enables flexible deployment and future extension, including integration with faster sampling methods (SnapFlow) and stronger base VLAs.

\end{document}
